\documentclass{article}
\begin{document}
\section*{Literal binary augmented naturality in ordered presentations}
The formal square is an equality of augmented native arrows for two
ordered dependent maps. The vertical maps are the original binary
Fubini isomorphisms. The horizontal native components are definitionally
the independently constructed occurrence-first and joint carrier
and whole-fibre maps, not replacement maps obtained by conjugation.
The occurrence-first route preserves the augmentation by identifying
its actual pair coefficients and endpoint maps with the independent
rectangular product of the two unary dependent computations. The
joint route's preservation follows from the previously proved
underlying native square and the vertical maps' invertibility as
augmented maps.

The source and target words are read in stable Cartesian order.
The comparison uses raw parent and child occurrence tags and never
deduplicates equal resource labels. Parent and continuation carrier
maps may be noninjective. Empty rows and globally empty words are
handled by the same structural proofs, using the supplied nullary
gathers and the specified whole maps from the additive unit, not
an assumption that the unit is initial. Syntactic zero subtrees
are retained until interpreted by the chosen gather.

The additive operation is supplied by a symmetric monoidal category.
The two multiplicative factor functors have explicit lax braided
structures whose nullary and binary maps agree with the declared
gathers. Multiplicative distributivity and its naturality,
the invertible middle and its compatibility, and the unit-middle
equation are inputs. The convenience theorem derives its unary
premises from the actual parent and continuation augmentation
laws and distinct addresses in each ordered presentation.

The binary augmented result varies over the specified ordered
inputs and their carrier-changing arrows. It does not yet prove
that the dependent or binary constructions preserve identities
or composition, does not identify the ordered square with an
order-independent finite-family natural transformation, and does
not establish the full all-arity dependent comparison in the source.
\end{document}
